\documentclass[fleqn,a4j,dvipdfmx,uplatex,twocolumn]{jsarticle}
\setlength{\columnseprule}{0.4pt}
\usepackage{amsmath,amssymb}
\usepackage{bm}
\usepackage{braket}
\usepackage[usenames]{color}
\usepackage[hiresbb]{graphicx}
\usepackage{enumerate}
\usepackage{prettyref}
\usepackage{polynom}
 \let\Ref\prettyref
 \newrefformat{eq}{\textbf{(\ref{#1})}}
 \newrefformat{sec}{第\ref{#1}節}
 \newrefformat{tab}{\textbf{表\ref{#1}}}
 \newrefformat{fig}{\textbf{図\ref{#1}}}

 \renewcommand{\Re}{\operatorname{Re}}
\renewcommand{\Im}{\operatorname{Im}}

\def\proof{\textbf{証明}　}
\def\qed{{\rule{0.5zw}{1zh}}}
\def\E{\mathrm{e}}
\def\D{\mathrm {d}}
\def\I{\mathrm {i}}
\def\LHS{(\mathrm{LHS})}
\def\RHS{(\mathrm{RHS})}
\def\hc{\mathrm{h.c.}}
\def\cc{\mathrm{c.c.}}
\def\ret{\notag\\&\qquad}
\def\nr{\notag\\}
\DeclareMathOperator{\Arg}{Arg}
\def\for#1{\quad\mathrm{for\ \ }#1}
\DeclareMathOperator{\sgn}{sgn}
\DeclareMathOperator{\Tr}{Tr}
\def\AAA{\mathaccent 23\mathrm{A}}
\def\abs#1{{\left\rvert #1 \right\lvert}}

\begin{document}
\begin{flushleft}
  新 基礎数学　改訂版，大日本図書
\end{flushleft}
\begin{center}
  \textbf{第1章 式の計算}\\
  \textbf{1　整式}\\
  \textbf{BASIC}\\
\end{center}

\paragraph{1 (1)}
\begin{align*}
  & 3 x^{2}-4 x+x^{2}-2+x+5 \\
  &= 4 x^{2}-3 x+3
\end{align*}

\paragraph{1 (2)}
\begin{align*}
  & 2+2 x^{2}-7 x+x^{2}+2 x-3+4 x \\
  &= 3 x^{2}-x-1
\end{align*}

\paragraph{2 (1)}
\begin{align*}
  A+B &= \left(2 x^{2}+5 x-3\right)+\left(3 x^{2}-3 x-3\right) \\
  &= 5 x^{2}+2 x-6 \\
  A-B &= \left(2 x^{2}+5 x-3\right)-\left(3 x^{2}-3 x-3\right) \\
  &= -x^{2}+8 x
\end{align*}

\paragraph{2 (2)}
\begin{align*}
  A+B &= \left(3 x^{3}-2 x+1\right)+\left(2 x^{3}-x^{2}+5 x\right) \\
  &= 5 x^{3}-x^{2}+3 x+1 \\
  A-B &= \left(3 x^{3}-2 x+1\right)-\left(2 x^{3}-x^{2}+5 x\right) \\
  &= x^{3}+x^{2}-7 x+1
\end{align*}

\paragraph{3 (1)}
\begin{align*}
  & 2 x^{2}-4 x y+3 y^{2}+3 x y+3 x^{2}-2 x-5 \\
  &= 5 x^{2}-x y-2 x+3 y^{2}-5 \\
  &= 5 x^{2}-(y+2) x+3 y^{2}-5
\end{align*}

\paragraph{3 (2)}
\begin{align*}
  & a^{3}-2 a^{2} x+a x^{2}+x^{3}-3 a x^{2}+a^{2} x \\
  &= x^{3}-2 a x^{2}-a^{2} x+a^{3}
\end{align*}

\paragraph{4 (1)}
\begin{align*}
  A+B &= \left(2 x^{3}+4 a x^{2}-3 a^{3}\right)+ \\
  & \left(3 a x^{2}-2 a^{2} x+2 a^{3}\right) \\
  &= 2 x^{3}+7 a x^{2}-2 a^{2} x-a^{3} \\
  A-B &= \left(2 x^{3}+4 a x^{2}-3 a^{3}\right)- \\
  & \left(3 a x^{2}-2 a^{2} x+2 a^{3}\right) \\
  &= 2 x^{3}+a x^{2}+2 a^{2} x-5 a^{3}
\end{align*}

\paragraph{4 (2)}
\begin{align*}
  A+B &= \left(x^{3}+2 x^{2} y+5 x y-y^{2}\right)+ \\
  & \left(3 x^{2} y-x y+2\right) \\
  &= x^{3}+5 x^{2} y+4 x y-y^{2}+2 \\
  &= -y^{2}+\left(5 x^{2}+4 x\right) y+x^{3}+2\\
  A-B &= (x^{3}+2 x^{2} y+5 x y-y^{2})-(3 x^{2} y-x y+2) \\
  &= x^{3}-x^{2} y+6 x y-y^{2}-2 \\
  &= -y^{2}-(x^{2}-6 x) y+x^{3}-2
\end{align*}

\paragraph{5 (1)}
\begin{align*}
  (-3 x)^{3} = -27 x^{3}
\end{align*}

\paragraph{5 (2)}
\begin{align*}
  (-a^{2})^{3} = -a^{6}
\end{align*}

\paragraph{5 (3)}
\begin{align*}
  (-a b^{3})^{2}(2 a^{2} b)^{3} &= a^{2} b^{6} \cdot 8 \cdot a^{6} \cdot b^{3} \\
  &= 8 a^{8} b^{9}
\end{align*}

\paragraph{5 (4)}
\begin{align*}
  (x+4)(x^{2}+3 x-5)= x^{3}+7 x^{2}+7 x-20
\end{align*}

\paragraph{6 (1)}
\begin{align*}
  (x+2 y)^{2} = x^{2}+4 x y+4 y^{2}
\end{align*}

\paragraph{6 (2)}
\begin{align*}
  (3 a-b)^{2} = 9 a^{2}-6 a b+b^{2}
\end{align*}

\paragraph{6 (3)}
\begin{align*}
  (3 x+2 y)(3 x-7 y)= 9 x^{2}-49 y^{2}
\end{align*}

\paragraph{6 (4)}
\begin{align*}
  (x+3 a)(x-2 a)= x^{2}+a x-6 a^{2}
\end{align*}

\paragraph{6 (5)}
\begin{align*}
  (2 x+3)(x+4)= 2 x^{2}+11 x+12
\end{align*}

\paragraph{6 (6)}
\begin{align*}
  (3 x-2 y)(2 x+3 y)= 6 x^{2}+5 x y-6 y^{2}
\end{align*}

\paragraph{6 (7)}
\begin{align*}
  (2 a+b)^{3} = 8 a^{3}+12 a^{2} b+6 a b^{2}+b^{3}
\end{align*}

\paragraph{6 (8)}
\begin{align*}
  (3 x-2 y)^{3} = 27 x^{3}-54 x^{2} y+36 x y^{2}-8 y^{3}
\end{align*}

\paragraph{7 (1)}
\begin{align*}
  (x+y+1)^{2} = x^{2}+y^{2}+1+2xy+2x+2y
\end{align*}

\paragraph{7 (2)}
\begin{align*}
  (x-2y+3)^{2} = x^{2}+4y^{2}+9-4xy+6x-12y
\end{align*}

\paragraph{7 (3)}
\begin{align*}
  (a+2)(a^{2}-2a+4) = a^{3}+8
\end{align*}

\paragraph{7 (4)}
\begin{align*}
  (3x-1)(9x^{2}+3x+1) = 27x^{3}-1
\end{align*}

\paragraph{8 (1)}
\begin{align*}
  & (2x+y+3)(2x+y+5) \\
  &= (2x+y)^{2}+8(2x+y)+15 \\
  &= 4x^{2}+4xy+y^{2}+16x+8y+15
\end{align*}

\paragraph{8 (2)}
\begin{align*}
  &  (a^{2}+a-1)(a^{2}-a+1) \\
  &= (a^{2})^{2}-(a-1)^{2} \\
  &= a^{4}-(a^{2}-2a+1) \\
  &= a^{4}-a^{2}+2a-1
\end{align*}

\paragraph{9 (1)}
\begin{align*}
  4a^{3}-9ab^{2} &= a(4a^{2}-9b^{2}) \\
  &= a(2a+3b)(2a-3b)
\end{align*}

\paragraph{9 (2)}
\begin{align*}
  ab-b-ac+c &= a(b-c)-(b-c) \\
  &= (b-c)(a-1)
\end{align*}

\paragraph{9 (3)}
\begin{align*}
  8a^{3}+27 = (2a+3)(4a^{2}-6a+9)
\end{align*}

\paragraph{9 (4)}
\begin{align*}
  x^{2}+2xy+y^{2}-z^{2} &= (x+y)^{2}-z^{2} \\
  &= (x+y+z)(x+y-z)
\end{align*}

\paragraph{10 (1)}
\begin{align*}
  x^{2}-5 x+6 = (x-2)(x-3)
\end{align*}

\paragraph{10 (2)}
\begin{align*}
  x^{2}-7 x-30 = (x-10)(x+3)
\end{align*}

\paragraph{11 (1)}
\begin{align*}
  6 x^{2}+17 x+5 = (2 x+5)(3 x+1)
\end{align*}

\paragraph{11 (2)}
\begin{align*}
  3 x^{2}+4 x-4 = (3 x-2)(x+2)
\end{align*}

\paragraph{12 (1)}
\begin{align*}
  x^{4}-16 &= (x^{2}+4)(x^{2}-4) \\
  &= (x^{2}+4)(x+2)(x-2)
\end{align*}

\paragraph{12 (2)}
\begin{align*}
  & (x-y)^{2}+2(x-y)-15 \\
  &= \{(x-y)+5\}\{(x-y)-3\} \\
  &= (x-y+5)(x-y-3)
\end{align*}

\paragraph{12 (3)}
\begin{align*}
  & x^{2}+2 x y+y^{2}-3 x-3 y+2 \\
  &= x^{2}+(2 y-3) x+(y-1)(y-2) \\
  &= \{x+(y-1)\}\{x+(y-2)\} \\
  &= (x+y-1)(x+y-2)
\end{align*}

\paragraph{12 (4)}
\begin{align*}
  & 2 x^{2}+5 x y+3 y^{2}-3 x-5 y-2 \\
  &= 2 x^{2}+(5 y-3) x+(3 y+1)(y-2) \\
  &= \{2 x+(3 y+1)\}\{x+(y-2)\} \\
  &= (2 x+3 y+1)(x+y-2)
\end{align*}

\paragraph{13 (1)}
\begin{align*}
  &\text{商 } x+5, \text{ 余り } 9 \\
  &\begin{array}{r}
    x + 5 \vspace{1pt} \\[-3pt]
    x - 2 \ \overline{) \ x^2 + 3x - 1} \\
    \underline{x^2 - 2x \phantom{-1}} \\
    5x - 1 \\
    \underline{5x - 10} \\
    9
  \end{array}
\end{align*}

\paragraph{13 (2)}
\begin{align*}
&\text{商 } 4x^2 - 13x + 16, \quad \text{余り } -16 \\
  &\begin{array}{r}
    4x^2 - 13x + 16 \vspace{1pt} \\[-3pt] % ここに商を記述
    x + 1 \ \overline{) \ 4x^3 - 9x^2 + 3x \phantom{+00}} \\
    \underline{4x^3 + 4x^2 \phantom{-000000}} \\
    -13x^2 + 3x \phantom{00} \\
    \underline{-13x^2 - 13x \phantom{00}} \\
    16x \phantom{00} \\
    \underline{16x + 16} \\
    -16
  \end{array}
\end{align*}

\paragraph{13 (3)}
\begin{align*}
&\text{商 } 2x + 1, \quad \text{余り } -9x + 3 \\
  &\begin{array}{r}
    2x + 1 \vspace{1pt} \\[-3pt]
    x^2 + x + 2 \ \overline{) \ 2x^3 + 3x^2 - 4x + 5} \\
    \underline{2x^3 + 2x^2 + 4x \phantom{+0}} \\
    x^2 - 8x + 5 \\
    \underline{x^2 + \phantom{0}x + 2} \\
    -9x + 3
  \end{array}
\end{align*}

\paragraph{14}
\begin{align*}
  &\text{ある整式を } A \text{ とすると} \nr
  A &= (2x+3)(3x^2-1)+5 \nr
  &= 6x^3+9x^2-2x-3+5 \nr
  &= 6x^3+9x^2-2x+2
\end{align*}

\paragraph{15 (1)}
\begin{align*}
  &\text{最大公約数 } bc \nr
  &\text{最小公倍数 } ab^2c^3d
\end{align*}

\paragraph{15 (2)}
\begin{align*}
  &\text{最大公約数 } x+2 \nr
  &\text{最小公倍数 } (x+2)(x-1)(x-2)
\end{align*}

\paragraph{15 (3)}
\begin{align*}
  &\text{最大公約数 } x+1 \nr
  &\text{最小公倍数 } (x+1)^2(x-1)(x^2-x+1)
\end{align*}

\paragraph{16 （1）}
\begin{align*}
  & 2 P(x)-Q(x) \\
  & =2\left(x^{3}-x^{2}-5 x+2\right)- \\
  & \left(2 x^{3}+6 x^{2}+3 x-1\right) \\
  & =-8 x^{2}-13 x+5
\end{align*}

\paragraph{16 （2）}
\begin{align*}
P(2) & =8-4-5 \cdot 2+2 \\
& =4-10+2 \\
& =-4
\end{align*}

\paragraph{16 （3）}
\begin{align*}
Q(-a)=-2 a^{3}+6 a^{2}-3 a-1
\end{align*}

\paragraph{17 (1)}
\begin{align*}
  A(2) &= 8 - 10 + 3 \\
       &= 1
\end{align*}

\paragraph{17 (2)}
\begin{align*}
  A(-3) &= -27 + 9 + 9 + 6 \\
        &= -3
\end{align*}

\paragraph{18}
\begin{align*}
  \intertext{$x = -\frac{1}{2}$ を代入すると}
  &2 \cdot \left( -\frac{1}{8} \right) - \frac{1}{4} - \frac{3}{2} + 5 \\
  &= -\frac{1}{4} - \frac{1}{4} + \frac{7}{2} \\
  &= -\frac{1}{2} + \frac{7}{2} \\
  &= 3
\end{align*}

\paragraph{19}
\begin{align*}
  P(-1) &= -1 + 4 - 1 - 6 = -4 \\
  P(-2) &= -8 + 16 - 2 - 6 = 0 \\
  P(-3) &= -27 + 36 - 3 - 6 = 0 \\
  \intertext{よって，$x+2, x+3$ で割り切れる}
\end{align*}

\paragraph{20}
\begin{align*}
  \intertext{$x = 1$ を代入した値が $0$ となるので}
  2 + k + 3 - 10 &= 0 \\
  k &= 5
\end{align*}

\paragraph{21 (1)}
\begin{align*}
  &x^3 - 7x + 6 \\
  &= (x - 1)(x^2 + x - 6) \\
  &= (x - 1)(x + 3)(x - 2)
\end{align*}

\paragraph{21 (2)}
\begin{align*}
  &x^3 + 4x^2 + 5x + 2 \\
  &= (x + 1)(x^2 + 3x + 2) \\
  &= (x + 1)^2(x + 2)
\end{align*}

\paragraph{21 (3)}
\begin{align*}
  &2x^3 + 3x^2 - 11x - 6 \\
  &= (x - 2)(2x^2 + 7x + 3) \\
  &= (x - 2)(2x + 1)(x + 3)
\end{align*}

\paragraph{21 (4)}
\begin{align*}
  &x^4 + 5x^3 + 5x^2 - 5x - 6 \\
  &= (x - 1)(x^3 + 6x^2 + 11x + 6) \\
  &= (x - 1)(x + 1)(x^2 + 5x + 6) \\
  &= (x - 1)(x + 1)(x + 2)(x + 3)
\end{align*}

\end{document}
